\section{System and design-space model}
\label{sec:model}

\subsection{Architecture template}
We consider the four-tier template of Fig.~\ref{fig:arch}. A population of devices
generates requests; each request traverses three application stages. The
\emph{pre-processing} stage $\tau_1$ runs either on the edge tier or on the gateway
tier, the \emph{aggregation} stage $\tau_2$ always runs on the gateway tier, and the
\emph{analytics} stage $\tau_3$ runs either on the edge tier or in the cloud. Which of
these variants is instantiated is a design decision (the placement policy).

\begin{figure}[t]
\centering
\begin{tikzpicture}[
  node distance=6mm and 9mm, font=\small,
  box/.style={draw, rounded corners=2pt, minimum height=7mm, minimum width=15mm,
              align=center, inner sep=2pt},
  lbl/.style={font=\scriptsize, midway, above, inner sep=1pt},
  arr/.style={-{Latex[length=2mm]}, thick}]
\node[box, fill=black!5] (dev) {devices\\[-2pt]\scriptsize $n_d$, rate $\lambda$};
\node[box, right=of dev, fill=edgec!15, draw=edgec] (edge)
     {edge tier\\[-2pt]\scriptsize $n_e \times$ class};
\node[box, right=of edge, fill=gwc!15, draw=gwc] (gw)
     {gateway tier\\[-2pt]\scriptsize $n_g \times$ class};
\node[box, right=of gw, fill=cloudc!12, draw=cloudc] (cl)
     {cloud\\[-2pt]\scriptsize class};
\draw[arr] (dev) -- node[lbl] {access} (edge);
\draw[arr] (edge) -- node[lbl] {access} (gw);
\draw[arr] (gw) -- node[lbl] {WAN} (cl);
\draw[arr, dashed] (dev.south) to[out=-25, in=-155] node[lbl, below]
     {\scriptsize cloud-first policy} (gw.south);
\draw[arr, dashed] (gw.north) to[out=155, in=25] node[lbl] {\scriptsize edge-first policy}
     (edge.north);
\node[below=1mm of edge, font=\scriptsize, text=edgec] {replication $r$};
\end{tikzpicture}
\caption{Architecture template. Solid arrows show the default request path; dashed
arrows show the alternative stage placements selected by the placement policy.}
\label{fig:arch}
\end{figure}

A compute node class is a tuple $(C, M, P^{\mathrm{idle}}, P^{\mathrm{peak}}, c)$ of
compute capacity [MIPS], memory [MB], idle and peak power [W] and normalised cost per
node. A link class is a tuple $(B, \tau^{\mathrm{prop}}, e, c)$ of bandwidth [Mb/s],
propagation delay [s], transfer energy [J/MB] and cost per attached node. The
gateway-to-cloud connection is fixed infrastructure ($B=200$\,Mb/s,
$\tau^{\mathrm{prop}}=20$\,ms) and is not a design variable. Table~\ref{tab:resources}
lists the concrete classes used in the experiments.

\begin{table}[t]
\centering
\caption{Resource classes. Compute classes: capacity $C$ [MIPS], memory $M$ [MB], idle
and peak power [W], cost [units]. Link classes: bandwidth $B$, propagation delay $\tau$,
transfer energy $e$, cost.}
\label{tab:resources}
\scriptsize
\begin{tabular}{lrrrrr}
\toprule
class & $C$ / $B$ & $M$ / $\tau$ & $P^{\mathrm{idle}}$ / $e$ & $P^{\mathrm{peak}}$ & cost \\
\midrule
\multicolumn{6}{l}{\emph{compute node classes}} \\
edge-low & 2000 & 2048 & 4 & 12 & 1 \\
edge-med & 6000 & 8192 & 8 & 30 & 2.4 \\
edge-high & 15000 & 16384 & 15 & 65 & 5.2 \\
gw-low & 4000 & 4096 & 6 & 20 & 1.8 \\
gw-high & 12000 & 16384 & 12 & 48 & 4 \\
cloud-med & 40000 & 131072 & 0 & 180 & 6 \\
cloud-high & 100000 & 262144 & 0 & 380 & 12 \\
\midrule \multicolumn{6}{l}{\emph{link classes} (name, $B$ [Mb/s], $\tau$ [ms], $e$ [J/MB], cost)} \\
constrained & 10 & 4 & 0.3 & 0.2 & -- \\
standard & 100 & 2 & 0.12 & 0.6 & -- \\
fast & 1000 & 1 & 0.05 & 1.4 & -- \\
\bottomrule
\end{tabular}

\end{table}

\subsection{Workload}
A workload is $(n_d, \lambda, \{\tau_i\}_{i=1}^{3}, D, s^{\mathrm{sync}})$ where $n_d$
devices each emit requests at rate $\lambda$, giving a system arrival rate
$\Lambda = n_d\lambda$; stage $\tau_i=(w_i, m_i, s_i)$ has computational work $w_i$ [MI],
resident memory $m_i$ [MB] and inbound payload $s_i$ [MB]; $D$ is the end-to-end deadline
and $s^{\mathrm{sync}}$ the per-request replica synchronisation payload. Three workload
classes are used throughout (Table~\ref{tab:workloads}); they are qualitatively
different rather than numerically rescaled versions of one another: telemetry is
high-rate and small-payload, real-time control is deadline-dominated, and data-intensive
sensing is payload- and compute-dominated.

\begin{table}[t]
\centering
\caption{Workload classes. Stage values are given as
$\tau_1/\tau_2/\tau_3$.}
\label{tab:workloads}
\scriptsize
\begin{tabular}{lrrrrrrr}
\toprule
workload & $n_d$ & $\lambda$ & $\Lambda$ & $w_i$ [MI] & $s_i$ [kB] & $D$ [ms] & $C^{\max}$ \\
\midrule
telemetry & 200 & 2 & 400 & 8/4/20 & 20/10/10 & 120 & 60 \\
control & 60 & 10 & 600 & 5/3/9 & 5/4/3 & 30 & 60 \\
sensing & 40 & 1.5 & 60 & 120/30/260 & 1200/350/300 & 900 & 90 \\
\bottomrule
\end{tabular}

\end{table}

\subsection{Decision vector and design space}
An architecture is described by the vector
\begin{equation}
x=\big(n_e,\ \kappa_e,\ n_g,\ \kappa_g,\ \ell,\ \kappa_c,\ r,\ \pi\big),
\label{eq:decision}
\end{equation}
with $n_e, n_g$ the number of edge nodes and gateways, $\kappa_e,\kappa_g,\kappa_c$ the
edge, gateway and cloud classes, $\ell$ the access link class, $r$ the replication
factor of stage $\tau_1$ and $\pi$ the placement policy. The design space is the
Cartesian product $\X=D_1\times\dots\times D_k$ with $k=8$, and
$|\X|=\prod_i |D_i|$. Table~\ref{tab:variables} lists the domains of the four scales
used in the experiments.

\begin{table}[t]
\centering
\caption{Design variables and domain sizes of the four design-space scales.}
\label{tab:variables}
\scriptsize
\begin{tabular}{llrrrr}
\toprule
variable & meaning & S1 & S2 & S3 & S4 \\
\midrule
$\mathit{n\_e}$ & number of edge nodes & 3 & 6 & 24 & 80 \\
$\mathit{edge\_cls}$ & edge node class & 3 & 3 & 3 & 3 \\
$\mathit{n\_g}$ & number of gateways & 2 & 3 & 5 & 8 \\
$\mathit{gw\_cls}$ & gateway class & 2 & 2 & 2 & 2 \\
$\mathit{link\_cls}$ & access link class & 3 & 3 & 3 & 3 \\
$\mathit{cloud\_cls}$ & cloud instance class & 2 & 2 & 2 & 2 \\
$\mathit{rep}$ & replication factor of stage 1 & 2 & 3 & 4 & 5 \\
$\mathit{policy}$ & placement policy & 3 & 3 & 3 & 3 \\
\midrule
\multicolumn{2}{l}{\emph{design-space size} $|\mathcal{X}|$} & 1\,296 & 5\,832 & 51\,840 & 345\,600 \\
\bottomrule
\end{tabular}

\end{table}

\subsection{Evaluation model}
Let $t(i)$ denote the tier executing stage $i$ under policy $\pi$, $N_t$ the number of
nodes of tier $t$ and $C_t$ the per-node capacity of its class. The offered load and
utilisation of a tier are
\begin{equation}
\Theta_t=\sum_{i:\,t(i)=t}\Lambda\,w_i\,\mu_i,\qquad
\rho_t=\frac{\Theta_t}{N_t C_t},\qquad
\mu_i=\begin{cases} r & i=1\\ 1 & \text{otherwise,}\end{cases}
\label{eq:rho}
\end{equation}
that is, replication multiplies the load of the ingest stage but not of the others.

\paragraph{Latency}
We use a stage-wise M/D/1-inspired waiting approximation
\cite{kleinrock1975queueing,harchol2013performance}: each stage retains its
per-node service time $w_i/C_{t(i)}$, while the waiting multiplier uses the
average tier utilisation $\rho_t$ based on aggregate capacity $N_tC_t$.
This is not an exact M/D/1 queue with service capacity $N_tC_t$: such a queue
would also divide the service time by $N_t$. Nor is it an exact multi-server
model. Stages sharing a tier have different service demands, and routing,
replication and successive visits can violate the Poisson-arrival assumption.
The expression is therefore a tractable surrogate whose prediction error must
be assessed empirically; it is not a proven lower bound on deployment latency.
The simulator and testbed implement per-node servers and quantify this gap.
With
$\mathcal{H}(x)$ the set of communication hops implied by the policy and the replication
factor,
\begin{equation}
L(x)=\underbrace{\sum_{i=1}^{3}\frac{w_i}{C_{t(i)}}}_{\text{service}}
+\underbrace{\sum_{i=1}^{3}\frac{w_i}{C_{t(i)}}\cdot\frac{\rho_{t(i)}}{2\,(1-\rho_{t(i)})}}_{\text{waiting}}
+\underbrace{\sum_{(s,\ell)\in\mathcal{H}(x)}\Big(\frac{8s}{B_\ell}+\tau^{\mathrm{prop}}_\ell\Big)}_{\text{communication}} .
\label{eq:latency}
\end{equation}
The middle term is what makes the objective \emph{contention-dependent}: it is a
function of $\rho_t$, which by \eqref{eq:rho} depends jointly on $n_t$, the class of the
tier, the replication factor and the policy. If $\rho_t\ge1$ for any used tier the
architecture is discarded as unstable.

\paragraph{Energy, cost and network volume}
Over an accounting horizon $T$,
\begin{align}
E(x) &= \sum_{t\in\{e,g\}} N_t\Big(P^{\mathrm{idle}}_t+\rho_t\big(P^{\mathrm{peak}}_t-P^{\mathrm{idle}}_t\big)\Big)T
      + \rho_c P^{\mathrm{peak}}_c T
      + \sum_{(s,\ell)\in\mathcal{H}(x)} s\Lambda e_\ell T, \label{eq:energy}\\
C(x) &= n_e\big(c_{\kappa_e}+c_\ell\big)+n_g\big(c_{\kappa_g}+c_\ell\big)+c_{\kappa_c},
\label{eq:cost}\\
D(x) &= \Lambda \sum_{(s,\ell)\in\mathcal{H}(x)} s .
\label{eq:netvol}
\end{align}
The objective vector is $F(x)=\big(L(x),E(x),C(x),D(x)\big)$, all components minimised.
Energy is the total consumption over the fixed accounting horizon $T=1$\,h and is
reported in kWh; the cost coefficients of Table~\ref{tab:resources} are synthetic
normalised values encoding the relative price tiers of the resource classes, and no
monetary interpretation is claimed. Neither the energy nor the cost model is validated
against measurements in this paper: they are used for comparative exploration only.

\paragraph{Feasibility}
An architecture is feasible when all of the following hold: $\rho_t\le\rho^{\max}$ for
every used tier; the resident memory of the stages mapped to a tier fits in the memory of
its node class; the aggregate access-link traffic does not exceed the total bandwidth of
the attached nodes; $L(x)\le D$; $C(x)\le C^{\max}$; and $E(x)\le E^{\max}$. We write
$\F$ for the set of feasible architectures.
